\documentclass[11pt]{article}
\usepackage[margin=1in]{geometry}
\usepackage{setspace}
\usepackage{hyperref}
\setstretch{1.15}

\begin{document}
\thispagestyle{empty}

\begin{flushright}
September 6, 2026
\end{flushright}

\vspace{1em}
\noindent
\textbf{To the Editor-in-Chief,}\\
\textit{Informatics in Medicine Unlocked}\\
Elsevier

\vspace{1.5em}
\noindent
\textbf{Subject: Submission of Original Research Article}

\vspace{1.5em}
\noindent
Dear Editor,

\vspace{1em}
\noindent
Please find enclosed our manuscript entitled \textbf{``Geodesic Diagnostic Trajectories (GDT): A Differential Geometry Framework for Adaptive Decision Support and Operational Optimization in EHR Systems''} by Hemanth Manchabale Papachappa and Aniriuddha Ganguly, which we are submitting for consideration as an Original Research article in \textit{Informatics in Medicine Unlocked}.

\vspace{1em}
\noindent
During electronic health record (EHR) chart review, clinicians frequently pivot between disjoint clinical topics (e.g., from sepsis markers to vasopressor dosing). Standard search systems fail to adapt to these abrupt shifts, suffering from ``latent-context pollution'' where stale query history actively degrades the retrieval of new, critical information. 

\vspace{1em}
\noindent
To address this challenge, we introduce \textbf{Geodesic Diagnostic Trajectories (GDT)}, a novel differential geometry framework that models evolving clinician intent as a trajectory on a Riemannian manifold. By mathematically formalizing topic pivots using a geodesic shift ratio, GDT elegantly suppresses irrelevant historical context with provable interference bounds. Coupled with a JEPA-style predictive prefetch engine, our approach reduces information retrieval latency by nearly 27\% compared to standard baselines. Using M/M/c queueing theory, we project this translates to an estimated time savings of 4.1 minutes per clinician during a peak-acuity 4-hour shift. 

\vspace{1em}
\noindent
We believe this manuscript is an excellent fit for the aims and scope of \textit{Informatics in Medicine Unlocked}. It directly bridges advanced computational AI methods with highly practical clinical applications, ultimately focusing on optimizing clinical decision support and EHR operational efficiency. Furthermore, our rigorous evaluation on multiple critical care databases (MIMIC-III/IV, eICU, and HiRID) provides actionable insights into the limitations of current sparse lexical retrieval systems, highlighting a clear pathway for modernizing clinical search tools.

\vspace{1em}
\noindent
We confirm that this manuscript is original, has not been published previously, and is not currently under consideration for publication elsewhere. All authors have reviewed and approved the final manuscript. We declare no competing financial interests or personal relationships that could have appeared to influence this work. To ensure reproducibility, our evaluation pipelines and code have been fully open-sourced.

\vspace{1em}
\noindent
Thank you for your time and consideration of our work. We look forward to your response.

\vspace{2em}
\noindent
Sincerely,\\
\vspace{1em}\\
\textbf{Hemanth Manchabale Papachappa}\\
% Uncomment and add your affiliation/contact details below if desired:
% Department of ... \\
% University / Institution \\
% Email: hemanth@example.com

\end{document}
